% Section 8 — Conclusion
\section{Conclusion}
\label{sec:conclusion}

\subsection{Answer to the research question}
\label{sec:answer}
Effectiveness does not march in step with sophistication. Moving from
traditional methods to neural forecasting reliably improves \emph{forecast
accuracy} --- the LSTM tops both datasets --- but improves \emph{inventory
outcomes} only conditionally: decisively on sparse intermittent demand, not at
all on dense demand under the default cost structure. The journey up the model
ladder buys capability everywhere and value somewhere.

\subsection{Main evidence}
\label{sec:main-evidence}
M5: LSTM best forecaster (MASE 1.316) and cheapest stocker (152.83), first in
25 of 27 policies --- robust. Store: LSTM best forecaster (MASE 0.978) but
fourth-cheapest stocker (2,247.45) behind Moving Average (2,084.49), SES, and
Seasonal Naive --- fragile (9/7/6/5 policy split). Ninety-one paired tests
separate the significant from the important (M5 LSTM--TSB: $p\approx2\times
10^{-10}$ yet $d_z=-0.031$; Store LSTM--MA: $d_z=-0.29$). H1 confirmed, H2
confirmed-with-qualification, H3 rejected in monotone form, H4 rejected in
strong form.

\subsection{Practical takeaway}
\label{sec:takeaway}
If your shelves look like M5 --- sparse, intermittent --- invest in the
learner; the forecast win reaches the cash register and survives policy
changes. If your shelves look like Store demand --- dense, smooth --- run a
moving average against anything expensive under your own cost ratio before you
buy complexity; and whatever wins, re-check it when lead times or penalties
move. In all cases, demand an inventory-level evaluation before trusting an
accuracy leaderboard.

\subsection{Final research conclusion}
\label{sec:final-word}
Forecasting skill and decision value are related but distinct quantities,
joined by demand structure and cost asymmetry. This study measured both, mapped
where they align and where they break, stress-tested the map across 27
policies, and froze every number for audit. The shelf, not the scoreboard, is
where forecasting earns its keep --- and now there is a reproducible record of
exactly when the smartest predictor is, and is not, the cheapest stocker.
